\chapter{Scheduling}
Why is scheduling such a big deal? Just brute force it and call it a day, should be as simple as that, right? The answer to that is a resounding no!

The scheduling problem we are dealing with here is deceptively large, even for small Unconferences. Take a small example: three rooms, five time slots, and twenty-one potential sessions to schedule. There are $\binom{n}{k}$ combinations for which sessions will be scheduled, where n is the total number of sessions and k is the capacity of the schedule $\binom{21}{15}=54,264$. Then for each of these combinations there are $\frac{n!}{k_1! \cdot k_2! \cdot k_3! \cdot k_4 \cdot k_5}$ ways to group the sessions across the 5 timeslots, where n is the number of sessions to schedule and $k_1$, $k_2$, $k_3$, $k_4$, and $k_5$ are the capacities of the 5 timeslots. In this schedule that works out to $\frac{15!}{3! \cdot 3! \cdot 3! \cdot 3! \cdot 3!}=168,168,000$. This means the total number of possible schedules is $num\_combinations \cdot ways\_to\_group$, or $54,264 \cdot 168,168,000=9,125,468,352,000$ possibilities for our scheduler to create and score just to determine the best schedule. At this size the problem is already getting too large to brute force in a reasonable amount of time. 

\section{Search}
There are several search concepts to consider: informed search (search that leverages heuristics: functions that indicate how good or close to the goal state they are), uninformed search (does not leverage information to guide it toward the goal), complete search (guaranteed to eventually find the goal), systematic search (guaranteed to find the goal without re-visiting any states), stochastic search (randomized, so not guaranteed to ever find the goal), local search (a common and effective kind of stochastic search), and anytime property (able to stop at any point and return the best solution so far, with longer runtime leading to better results and eventually optimality).

The problem with brute force search is that it is an uninformed systematic search that explores areas that move it away from the goal. If brute force is not a reasonable solution, how do we go about scheduling? That depends on our goals for scheduling. If the goal is to determine whether a solution exists or to find the optimal solution, then a complete or systematic search would work well. However, for UnconfRS, the goal is to find a reasonably good solution quickly. This makes informed local search a good fit since it lends itself to the anytime property, allowing the scheduler to return the best solution found so far.
%
%
\section{Local Search}
The local search used in this project works by starting at an initial candidate solution and iteratively moving to neighboring states in the search space (the set of all unique solutions). A neighbor (a different candidate solution from a single small move in the previous state) for UnconfRS would be swapping a session with another, either already on the schedule or with one that is unassigned.

At each step, the neighborhood (the set of all neighbors from the current state) is evaluated using heuristic functions. A neighbor is then chosen to explore. Which neighbor is chosen depends on the local search algorithm, ranging from greedy (best improving) to stochastic. This continues until reaching a local minimum (candidate solution is better than all its neighbors). It is possible it is also a global minimum (candidate solution better than all others in the search space). However, there is no guarantee of finding a global minimum using local search.

Random restarts and noise moves are strategies that can be used to avoid getting trapped in a non-global minimum. With random restarts, the scheduling is run many times from different randomly initialized starting states. Noise moves are random moves that attempt to prevent the search from getting stuck in a local minimum.